\documentclass[11pt]{article}
\usepackage[margin=1in]{geometry}
\usepackage{amsmath, amssymb}
\usepackage{graphicx}
\usepackage{hyperref}

\title{\textbf{HW\#00: MLP function approximation [100 pts]}}
\author{Scientific Machine Learning}
\date{}
\begin{document}
\maketitle
\hrulefill

\section{Theoretical Foundations [25 pts]}

\subsection{The Universal Approximation Theorem (UAT)}
\begin{enumerate}
    \item[\textbf{a)}] In 2-3 lines, state the Universal Approximation Theorem for a single-hidden-layer feedforward neural network.
    \item[\textbf{b)}] Briefly discuss the limitations of the UAT. What does it guarantee, and what does it \textit{not} guarantee about network training, architecture, or efficiency?
\end{enumerate}

\subsection{Bernstein Polynomials and a ReLU Network}
The Bernstein polynomial of degree $n$ for $f$ on $[0,1]$ is
\[ B_n f(x) = \sum_{k=0}^{n} f\!\left(\tfrac{k}{n}\right) \binom{n}{k} x^k (1-x)^{n-k}. \]
Take $f_1(x) = \sin(\pi x)$ and $f_2(x) = |x - \tfrac12|$.
\begin{enumerate}
    \item[\textbf{a)}] For each function, plot $B_n f$ against $f$ for $n = 4, 16, 64, 256$ and report the largest gap $\max_x |f(x) - B_n f(x)|$ for each $n$.
    \item[\textbf{b)}] Repeat with a one-hidden-layer ReLU network of width $4, 16, 64, 256$, trained on $200$ evenly spaced samples of $f$.
    \item[\textbf{c)}] Which function does each representation handle better, and why?
\end{enumerate}

\subsection{Derivatives of Sigmoid and ReLU}
\begin{enumerate}
    \item[\textbf{a)}] For $\sigma(z) = (1 + e^{-z})^{-1}$, derive $\sigma'(z)$ by hand and show that $\sigma'(z) = \sigma(z)\,(1 - \sigma(z))$.
    \item[\textbf{b)}] For $\mathrm{ReLU}(z) = \max(0, z)$, write the derivative for $z < 0$ and for $z > 0$. What happens at $z = 0$?
    \item[\textbf{c)}] Plot $\sigma$, $\sigma'$, $\mathrm{ReLU}$, and $\mathrm{ReLU}'$ on $[-5, 5]$. Where is $\sigma'$ small, and what does that do to a gradient that passes through such a unit during backpropagation?
\end{enumerate}

\section{Computational Experiments: Function Approximation [25 pts]}

\subsection{Approximating a Discontinuous Function}
This problem explores the challenge of approximating non-smooth functions. Consider the Heaviside step function on the interval $x \in [0,1]$:
\[ H(x) = \begin{cases} 0 & x < 0.5 \\ 1 & x \geq 0.5 \end{cases} \]
\begin{enumerate}
    \item[\textbf{a)}] Using a single hidden layer MLP with 50 neurons and a `tanh` activation, train the network to approximate $H(x)$.
    \item[\textbf{b)}] Plot the learned function $\hat{u}(x)$ against the true function $H(x)$. Describe the behavior around the discontinuity at $x=0.5$. Does it resemble the Gibbs phenomenon?
    \item[\textbf{c)}] Experiment to improve the approximation by independently:
        \begin{itemize}
            \item Increasing the hidden layer width (e.g., to 200, then 500 neurons).
            \item Increasing the network depth (e.g., 3 hidden layers of 50 neurons each).
            \item Changing the activation function to ReLU.
            \item Try different loss functions
        \end{itemize}
    \item[\textbf{d)}] \textbf{Analysis:} Why does a `tanh`-based network struggle with discontinuities? Which modification produced the best result and why? Relate your findings to the fact that compositions of `tanh` are $C^\infty$, while ReLU networks are piecewise linear.
\end{enumerate}

\section{Compositional Learning [25 pts]}
This section explores how to solve complex problems by breaking them into simpler, modular components that can be learned separately and then combined.

\subsection{Learning and Composing Multi-Scale Functions}
Consider the function $h(x) = f_{\text{linear}}(x) + f_{\text{high}}(x)$ on the interval $x \in [0,1]$, where:
\begin{itemize}
    \item $f_{\text{linear}}(x) = 1 + 2x$ (a simple linear trend)
    \item $f_{\text{high}}(x) = 0.5 \cos(20\pi x)$ (a high-frequency component)
\end{itemize}

\subsubsection*{Part 1: Learning the High-Frequency Component}
To understand the challenge, first focus solely on approximating $f_{\text{high}}(x)$.
\begin{enumerate}
    \item[\textbf{a)}] Train MLPs of increasing depth (1, 2, 3, and 4 hidden layers) to learn $f_{\text{high}}(x)$.
    \item[\textbf{b)}] Plot the results on a 2x2 grid to visually demonstrate how approximation quality improves with depth. Explain this result in the context of spectral bias.
\end{enumerate}

\subsubsection*{Part 2: Composition by ``Stitching" Networks}
Now, we will tackle the full problem $h(x)$ using a ``divide and conquer" strategy.
\begin{enumerate}
    \item[\textbf{a)}] \textbf{Modular Training:}
        \begin{itemize}
            \item Train a simple, shallow network, $N_{\text{linear}}$, to learn $f_{\text{linear}}(x)$.
            \item Train a separate, deep network, $N_{\text{high}}$, to learn $f_{\text{high}}(x)$.
        \end{itemize}
    \item[\textbf{b)}] \textbf{Manual Composition:} The final stitched model is formed by adding the outputs of the two trained networks:
    \[ \hat{h}_{\text{stitched}}(x) = N_{\text{linear}}(x) + N_{\text{high}}(x) \]
    Implement this composition in code.
    \item[\textbf{c)}] \textbf{Comparison \& Analysis:} Train a single, deep end-to-end network to learn $h(x)$ directly. Create a side-by-side plot comparing the result of your stitched model against the end-to-end model. For clarity, zoom in on the interval $x \in [0.2, 0.8]$ for one of your plots. Discuss the results. Which approach is more effective and why?
\end{enumerate}



\section{Regularization [25 pts]}
\subsection{L1 vs. L2 Penalties}
Take $40$ evenly spaced points $x_i$ on $[0, 2\pi]$ with readings $y_i = \sin(x_i) + \varepsilon_i$, $\varepsilon_i \sim \mathcal{N}(0, 0.1^2)$.
Fit $\hat y(x) = \sum_{j=1}^{100} w_j \phi_j(x)$ with Gaussian bumps $\phi_j(x) = \exp\!\left(-(x - c_j)^2 / 2\ell^2\right)$, centers $c_j$ evenly spaced on $[0, 2\pi]$ and $\ell = 0.5$, by minimizing
\[ \frac{1}{40}\sum_{i=1}^{40} \bigl(\hat y(x_i) - y_i\bigr)^2 + \lambda R(\mathbf{w}), \qquad R(\mathbf{w}) = \sum_j w_j^2 \ \text{(L2)} \quad\text{or}\quad R(\mathbf{w}) = \sum_j |w_j| \ \text{(L1)}. \]
Solve the L1 problem by soft thresholding or with a library solver.
\begin{enumerate}
	\item[\textbf{a)}] Fit with no penalty. Plot the fit against $\sin(x)$ and the readings. What happens between the readings, and why?
	\item[\textbf{b)}] Fit with L2 and with L1 for $\lambda \in \{10^{-4}, 10^{-3}, 10^{-2}, 10^{-1}\}$. Report the RMSE against $\sin(x)$ for each.
	\item[\textbf{c)}] At $\lambda = 10^{-2}$, plot the $100$ weights of the L2 and L1 fits. How many are exactly zero in each, and why?
\end{enumerate}

\end{document}
